% GENERATED FILE. Edit canonical lessons, then run npm run generate:books.
% canonical-source-hash: fnv1a64:8a9c6d8c1e8a2486
% canonical-lessons: MW-C215-lapetno, MW-C215-kagad, MW-C215-potli, MW-C215-ochho, MW-C215-khutno, MW-R215-review-210-215, MW-R215-review-final-215

\chapter{Wrapping Up}
\label{ch:mw-wrapping-up}
\hlchaptermodality{215}

\begin{chapteropening}
\textbf{By the end of this chapter:} \emph{I can wrap what I bought and say what has run out.}

Complete the last lesson of chapter 215: I can wrap what I bought and say what has run out.
\end{chapteropening}

\section[\texorpdfstring{\mw{लपेटणो} (lapeṭṇo)}{लपेटणो (lapeṭṇo)}]{\mw{लपेटणो} (lapeṭṇo) — to wrap}

\label{lesson:MW-C215-lapetno}



\begin{quote}
Before the new one: say the Marwadi for to flip over, then the Marwadi for to turn up.
\end{quote}

\subsection*{You'll want to know: \mw{लपेटणो}}
\textbf{\mw{लपेटणो}} — \emph{lapeṭṇo} — \textquotedblleft{}to wrap\textquotedblright{}.

\begin{quote}
\mw{मां} \mw{रोटी} \mw{कपड़ा} \mw{मांय} \mw{लपेटै} \mw{है।}
\end{quote}

\begin{quote}
\emph{mā̃ roṭī kapṛā mā̃y lapeṭai hai.}
\end{quote}

\begin{quote}
Mother wraps the flatbread in a cloth.
\end{quote}

Past (I …): \mw{म्हैं} \mw{रोटी} \mw{लपेटी}, \emph{mhaĩ roṭī lapeṭī}.

\begin{cousinweb}
\textbf{\mw{लपेटणो}} (\emph{lapeṭṇo}) is inherited from a reconstructed \emph{*lappēṭṭ}, \textquotedblleft{}wrap\textquotedblright{}. CDIAL gives no Marwari form in this entry, but lists Hindi \emph{lapeṭnā}, \textquotedblleft{}to wrap, fold, wind\textquotedblright{} (CDIAL 10942).
\end{cousinweb}

\begin{grammarlens}[title={What you've built}]
One more word for this chapter.
\end{grammarlens}

\subsection*{Your turn}
\begin{itemize}
\raggedright
  \item \emph{Say it:} \emph{lapeṭṇo}
  \item \emph{Say it:} \emph{lapeṭṇo}, once more
  \item \emph{Say it:} \emph{mā̃ roṭī kapṛā mā̃y lapeṭai hai.} — Mother wraps the flatbread in a cloth.
  \item \emph{Recall:} say the Marwadi for to flip over, then the Marwadi for to turn up, then say \emph{lapeṭṇo} again
\end{itemize}

\begin{tcolorbox}[breakable,colback=teal!4,colframe=teal!35!black,title={Before you move on}]
What does \mw{लपेटणो} mean? (\textquotedblleft{}To wrap\textquotedblright{}.) Say it once more.
\end{tcolorbox}
\section[\texorpdfstring{\mw{कागद} (kāgad)}{कागद (kāgad)}]{\mw{कागद} (kāgad) — paper; a letter}

\label{lesson:MW-C215-kagad}



\begin{quote}
Before the new one: say the Marwadi for to turn up, then the Marwadi for to wrap.
\end{quote}

\subsection*{You'll want to know: \mw{कागद}}
\textbf{\mw{कागद}} — \emph{kāgad} — \textquotedblleft{}paper; a letter\textquotedblright{}.

\begin{quote}
\mw{कागद} \mw{मेज} \mw{माथै} \mw{है।}
\end{quote}

\begin{quote}
\emph{kāgad mej māthai hai.}
\end{quote}

\begin{quote}
The paper is on the table.
\end{quote}

A masculine noun.

\begin{cousinweb}
\textbf{\mw{कागद}} (\emph{kāgad}) is usually traced to Persian \emph{kāghaz}, \textquotedblleft{}paper\textquotedblright{}.
\end{cousinweb}

\begin{grammarlens}[title={What you've built}]
One more word for this chapter.
\end{grammarlens}

\subsection*{Your turn}
\begin{itemize}
\raggedright
  \item \emph{Say it:} \emph{kāgad}
  \item \emph{Say it:} \emph{kāgad}, once more
  \item \emph{Say it:} \emph{kāgad mej māthai hai.} — The paper is on the table.
  \item \emph{Recall:} say the Marwadi for to turn up, then the Marwadi for to wrap, then say \emph{kāgad} again
\end{itemize}

\begin{tcolorbox}[breakable,colback=teal!4,colframe=teal!35!black,title={Before you move on}]
What does \mw{कागद} mean? (\textquotedblleft{}Paper; a letter\textquotedblright{}.) Say it once more.
\end{tcolorbox}
\section[\texorpdfstring{\mw{पोटली} (poṭlī)}{पोटली (poṭlī)}]{\mw{पोटली} (poṭlī) — a small bundle}

\label{lesson:MW-C215-potli}



\begin{quote}
Before the new one: say the Marwadi for to wrap, then the Marwadi for paper.
\end{quote}

\subsection*{You'll want to know: \mw{पोटली}}
\textbf{\mw{पोटली}} — \emph{poṭlī} — \textquotedblleft{}a small bundle\textquotedblright{}.

\begin{quote}
\mw{पोटली} \mw{मांय} \mw{रोटी} \mw{है।}
\end{quote}

\begin{quote}
\emph{poṭlī mā̃y roṭī hai.}
\end{quote}

\begin{quote}
There is flatbread in the bundle.
\end{quote}

A feminine noun.

\begin{cousinweb}
\textbf{\mw{पोटली}} (\emph{poṭlī}) is inherited from a reconstructed \emph{*pōṭṭa}, \textquotedblleft{}bundle\textquotedblright{}. CDIAL gives no Marwari form in this entry, but lists Hindi \emph{poṭlī}, \textquotedblleft{}small bundle\textquotedblright{} (CDIAL 8396).
\end{cousinweb}

\begin{grammarlens}[title={What you've built}]
One more word for this chapter.
\end{grammarlens}

\subsection*{Your turn}
\begin{itemize}
\raggedright
  \item \emph{Say it:} \emph{poṭlī}
  \item \emph{Say it:} \emph{poṭlī}, once more
  \item \emph{Say it:} \emph{poṭlī mā̃y roṭī hai.} — There is flatbread in the bundle.
  \item \emph{Recall:} say the Marwadi for to wrap, then the Marwadi for paper, then say \emph{poṭlī} again
\end{itemize}

\begin{tcolorbox}[breakable,colback=teal!4,colframe=teal!35!black,title={Before you move on}]
What does \mw{पोटली} mean? (\textquotedblleft{}A small bundle\textquotedblright{}.) Say it once more.
\end{tcolorbox}
\section[\texorpdfstring{\mw{ओछो} (ochho)}{ओछो (ochho)}]{\mw{ओछो} (ochho) — too small; short}

\label{lesson:MW-C215-ochho}



\begin{quote}
Before the new one: say the Marwadi for paper, then the Marwadi for a small bundle.
\end{quote}

\subsection*{You'll want to know: \mw{ओछो}}
\textbf{\mw{ओछो}} — \emph{ochho} — \textquotedblleft{}too small; short\textquotedblright{}.

\begin{quote}
\mw{म्हारी} \mw{टोपी} \mw{ओछी} \mw{है।}
\end{quote}

\begin{quote}
\emph{mhārī ṭopī ochhī hai.}
\end{quote}

\begin{quote}
My cap is too small.
\end{quote}

\begin{cousinweb}
\textbf{\mw{ओछो}} (\emph{ochho}) is inherited from a reconstructed \emph{*ōccha}, \textquotedblleft{}small, thin, mean\textquotedblright{}. CDIAL gives no Marwari form in this entry, but lists Hindi \emph{ochā}, \textquotedblleft{}empty, small\textquotedblright{}, and Gujarati \emph{ochũ}, \textquotedblleft{}deficient\textquotedblright{} (CDIAL 2540).
\end{cousinweb}

\begin{grammarlens}[title={What you've built}]
One more word for this chapter.
\end{grammarlens}

\subsection*{Your turn}
\begin{itemize}
\raggedright
  \item \emph{Say it:} \emph{ochho}
  \item \emph{Say it:} \emph{ochho}, once more
  \item \emph{Say it:} \emph{mhārī ṭopī ochhī hai.} — My cap is too small.
  \item \emph{Recall:} say the Marwadi for paper, then the Marwadi for a small bundle, then say \emph{ochho} again
\end{itemize}

\begin{tcolorbox}[breakable,colback=teal!4,colframe=teal!35!black,title={Before you move on}]
What does \mw{ओछो} mean? (\textquotedblleft{}Too small; short\textquotedblright{}.) Say it once more.
\end{tcolorbox}
\section[\texorpdfstring{\mw{खुटणो} (khuṭṇo)}{खुटणो (khuṭṇo)}]{\mw{खुटणो} (khuṭṇo) — to run out}

\label{lesson:MW-C215-khutno}



\begin{quote}
Before the new one: say the Marwadi for a small bundle, then the Marwadi for too small.
\end{quote}

\subsection*{You'll want to know: \mw{खुटणो}}
\textbf{\mw{खुटणो}} — \emph{khuṭṇo} — \textquotedblleft{}to run out\textquotedblright{}.

\begin{quote}
\mw{घर} \mw{मांय} \mw{आटो} \mw{खुटै} \mw{है।}
\end{quote}

\begin{quote}
\emph{ghar mā̃y āṭo khuṭai hai.}
\end{quote}

\begin{quote}
We're running out of flour at home.
\end{quote}

Past (it …): \mw{वो} \mw{आटो} \mw{खुट्यो}, \emph{vo āṭo khuṭyo}.

\begin{cousinweb}
\textbf{\mw{खुटणो}} (\emph{khuṭṇo}) is inherited from a reconstructed \emph{*khuṭṭa}, \textquotedblleft{}defective\textquotedblright{}. CDIAL gives no Marwari form in this entry, but lists Gujarati \emph{khuṭvũ}, \textquotedblleft{}to be exhausted\textquotedblright{}, and Sindhi \emph{khuṭaṇu}, \textquotedblleft{}to fall short\textquotedblright{} (CDIAL 3895).
\end{cousinweb}

\begin{grammarlens}[title={What you've built}]
One more word for this chapter.
\end{grammarlens}

\subsection*{Your turn}
\begin{itemize}
\raggedright
  \item \emph{Say it:} \emph{khuṭṇo}
  \item \emph{Say it:} \emph{khuṭṇo}, once more
  \item \emph{Say it:} \emph{ghar mā̃y āṭo khuṭai hai.} — We're running out of flour at home.
  \item \emph{Recall:} say the Marwadi for a small bundle, then the Marwadi for too small, then say \emph{khuṭṇo} again
\end{itemize}

\begin{tcolorbox}[breakable,colback=teal!4,colframe=teal!35!black,title={Before you move on}]
What does \mw{खुटणो} mean? (\textquotedblleft{}To run out\textquotedblright{}.) Say it once more.
\end{tcolorbox}
\section[(dialogue)]{Looking back: from Choosing the Right One to Wrapping Up}

\label{lesson:MW-R215-review-210-215}



\begin{quote}
Say the Marwadi for to test and a sort.
\end{quote}

\subsection*{Your turn}
Say these aloud:

\begin{itemize}
\raggedright
  \item to test, a sort, a sample, fake, to pick out
  \item a size, to measure, roomy, tight, to fit
  \item plus a quarter, one and a half, two and a half, double, a chhatak
  \item the cost, to decrease, to lower, free of charge, to be sold
  \item to give back, alike, a small bag, to flip over, to turn up
  \item to wrap, paper, a small bundle, too small, to run out
\end{itemize}

\begin{tcolorbox}[breakable,colback=teal!4,colframe=teal!35!black,title={Before you move on}]
To test? (\textbf{\mw{परखणो}}, \emph{parakhṇo}.) A sort? (\textbf{\mw{भांत}}, \emph{bhā̃t}.) A sample? (\textbf{\mw{नमूनो}}, \emph{namūno}.) Fake? (\textbf{\mw{नकली}}, \emph{naklī}.) To pick out? (\textbf{\mw{चुणणो}}, \emph{chuṇṇo}.) A size? (\textbf{\mw{माप}}, \emph{māp}.) To measure? (\textbf{\mw{नापणो}}, \emph{nāpṇo}.) Roomy? (\textbf{\mw{मोकळो}}, \emph{mokḷo}.) Tight? (\textbf{\mw{तंग}}, \emph{taṅg}.) To fit? (\textbf{\mw{समावणो}}, \emph{samāvṇo}.) Plus a quarter? (\textbf{\mw{सवा}}, \emph{savā}.) One and a half? (\textbf{\mw{डेढ}}, \emph{ḍeḍh}.) Two and a half? (\textbf{\mw{अढाई}}, \emph{aḍhāī}.) Double? (\textbf{\mw{दूणो}}, \emph{dūṇo}.) A chhatak? (\textbf{\mw{छटांक}}, \emph{chhaṭā̃k}.) The cost? (\textbf{\mw{मोल}}, \emph{mol}.) To decrease? (\textbf{\mw{घटणो}}, \emph{ghaṭṇo}.) To lower? (\textbf{\mw{घटावणो}}, \emph{ghaṭāvṇo}.) Free of charge? (\textbf{\mw{मुफत}}, \emph{muphat}.) To be sold? (\textbf{\mw{बिकणो}}, \emph{bikṇo}.) To give back? (\textbf{\mw{फेरणो}}, \emph{pherṇo}.) Alike? (\textbf{\mw{सारीखो}}, \emph{sārīkho}.) A small bag? (\textbf{\mw{थैली}}, \emph{thailī}.) To flip over? (\textbf{\mw{पलटणो}}, \emph{palaṭṇo}.) To turn up? (\textbf{\mw{लाधणो}}, \emph{lādhṇo}.) To wrap? (\textbf{\mw{लपेटणो}}, \emph{lapeṭṇo}.) Paper? (\textbf{\mw{कागद}}, \emph{kāgad}.) A small bundle? (\textbf{\mw{पोटली}}, \emph{poṭlī}.) Too small? (\textbf{\mw{ओछो}}, \emph{ochho}.) To run out? (\textbf{\mw{खुटणो}}, \emph{khuṭṇo}.)
\end{tcolorbox}
\section[(dialogue)]{Wrapping Up — once more}

\label{lesson:MW-R215-review-final-215}



\begin{quote}
Say the Marwadi for to wrap and paper.
\end{quote}

\subsection*{Your turn}
\begin{itemize}
\raggedright
  \item \emph{Say it:} to wrap, paper, a small bundle, too small, to run out
\end{itemize}

\begin{tcolorbox}[breakable,colback=teal!4,colframe=teal!35!black,title={Before you move on}]
To wrap? (\textbf{\mw{लपेटणो}}, \emph{lapeṭṇo}.) Paper? (\textbf{\mw{कागद}}, \emph{kāgad}.) A small bundle? (\textbf{\mw{पोटली}}, \emph{poṭlī}.) Too small? (\textbf{\mw{ओछो}}, \emph{ochho}.) To run out? (\textbf{\mw{खुटणो}}, \emph{khuṭṇo}.)
\end{tcolorbox}
